\chapter{مخارج الآلة: كيف يتحكم الدماغ في العضلات}
\chaptermark{مخارج الآلة}
\label{sec:muscles}

\section{المخرج السريع}

كل ما تفعله الآلة في العالم بسرعة تفعله بالعضلات. الكلمة، والنظرة، والخطوة، والابتسامة كلها انقباضات عضلات. أما المخرج الثاني البطيء، أي كيمياء الجسم، فنتناوله في الفصل~\ref{sec:chemoutput}. في الشكل~\ref{fig:machine} كان المخرج سهمًا عنوانه ”العضلات“؛ فلننظر الآن إلى ما وراء هذا السهم.

الحلقة الأخيرة في السلسلة عصبونات \nt{ACET}، وهي نفسها التي كُتب عنها في الجدول~\ref{tab:types} ”المخرج إلى العضلة“. كل ليف عضلي يتلقى دخله من عصبون واحد بالضبط من هذه العصبونات، والعصبون الواحد يتحكم في مجموعة من الألياف، من بضع مئات إلى نحو ألفين~\cite{Feinstein1955mu}. وفي العضلات التي تحتاج إلى ضبط دقيق تكون الوحدات أصغر، وفي عضلات الساقين الكبيرة أكبر. ويسمى العصبون مع أليافه \emph{الوحدة الحركية}: إنها أصغر قطعة من العضلة يستطيع الدماغ تشغيلها وحدها.

مشبك العصبون على العضلة مبني على نحو مختلف تمامًا عن مشابك القشرة في الفصل~\ref{sec:code}. هناك كان معظم النبضات يضيع. أما هنا فكل نبضة تطلق من الأستيل كولين أضعاف ما يلزم لإطلاق الليف~\cite{WoodSlater2001}. إنه مشبك ذو \emph{هامش أمان}: نبضة واحدة من العصبون تعني انقباضة واحدة بالضبط للألياف. داخل الآلة يمكن تحمّل وصلات غير موثوقة وتصحيح الأخطاء بالفائض؛ أما في المخرج فالخطأ يكلف حركة، والموثوقية مشتراة في العتاد مباشرة.

\section{القوة: التردد وعدد الوحدات}

النبضة الواحدة تعطي انقباضة مفردة، أي نفضة: تتصاعد القوة خلال بضع عشرات من الملّي ثانية ثم تهبط. وتقريب جيد لها~\cite{Fuglevand1993}
\begin{equation}
h(t) \;=\; F_1\,\frac{t}{T_c}\,e^{\,1-t/T_c},
\label{eq:twitch}
\end{equation}
حيث $T_c$ زمن التصاعد، من رتبة $50\ms$، و$F_1$ قوة النفضة الواحدة. وإذا تقاربت النبضات تجمّعت النفضات:
\begin{empheq}[box=\keybox]{equation}
F(t) \;=\; \sum_k h\bigl(t-t_k\bigr) .
\label{eq:forcesum}
\end{empheq}
هذا هو مرشح التمرير المنخفض نفسه الذي حوّل النبضات إلى تركيز في الفصل~\ref{sec:serocomputer}، إلا أن المخرج الآن ليس تركيزًا بل قوة. العضلة محوِّل رقمي تماثلي من النبضات إلى القوة (الشكل~\ref{fig:tetanus}). في نموذجنا، عند خمس نبضات في الثانية لا تكاد النفضات تتداخل، فتقفز القوة بين $0.2F_1$ و$1.1F_1$؛ وعند خمس عشرة تبقى بين $1.8F_1$ و$2.2F_1$؛ وعند أربعين تصبح شبه مستوية، نحو $5.4F_1$. والعضلة الحقيقية تتشبع مع النبضات المتقاربة، فلا يصح الجمع الخطي~\eqref{eq:forcesum} إلا حتى بضع عشرات من النبضات في الثانية.

\begin{figure}[t]
\centering
\begin{tikzpicture}[>=Stealth,x=20cm,y=0.55cm]
\draw[->,gray!70!black] (0,0) -- (0.66,0) node[right,font=\small,black] {$t,\ \text{s}$};
\draw[->,gray!70!black] (0,0) -- (0,6.4) node[left,font=\small,black] {$F/F_1$};
\foreach \t in {0,0.2,0.4,0.6}{\draw[gray!60!black] (\t,-0.1) -- (\t,0.1); \node[below,font=\scriptsize] at (\t,0) {\t};}
\foreach \f in {1,3,5}{\draw[gray!60!black] (-0.004,\f) -- (0.004,\f); \node[left,font=\scriptsize] at (0,\f) {\f};}
\draw[thick,blue!60!black] plot file {figures/muscle_twitch_5.dat};
\draw[thick,green!45!black] plot file {figures/muscle_twitch_15.dat};
\draw[thick,red!70!black] plot file {figures/muscle_twitch_40.dat};
\node[font=\scriptsize,blue!60!black,anchor=west] at (0.605,0.9) {5 نبضات/ث};
\node[font=\scriptsize,green!45!black,anchor=west] at (0.605,2.1) {15 نبضة/ث};
\node[font=\scriptsize,red!70!black,anchor=west] at (0.605,5.4) {40 نبضة/ث};
\end{tikzpicture}
\caption{العضلة محوِّلًا رقميًا تماثليًا: القوة~\eqref{eq:forcesum} عند نبضات منتظمة من وحدة حركية واحدة، $T_c=50\ms$. النبضات المتباعدة تعطي نفضات منفصلة، والمتقاربة تندمج في قوة مستوية، كما تندمج النبضات المتقاربة في تركيز مستوٍ في الشكل~\ref{fig:lowpass}.}
\label{fig:tetanus}
\end{figure}

للدماغ مقبضان للقوة: تردد نبضات كل وحدة، وعدد الوحدات المشغَّلة. والمقبض الثاني مبني ببراعة كبيرة. الوحدات في العضلة مختلفة: أضعفها أضعف من أقواها مئة مرة. وحين تلزم قوة متزايدة تُشغَّل الوحدات \emph{بترتيب الحجم}، من الصغيرة إلى الكبيرة~\cite{Henneman1957}. لا أحد يحسب هذا الترتيب: العصبونات الصغيرة ببساطة أسهل استثارة بالدخل المشترك نفسه، فترتيب التشغيل يحدده العتاد نفسه.

ماذا يعطي ذلك؟ لنفترض أن كل وحدة تالية في الترتيب أقوى من سابقتها بـ$q$ مرة، حيث $q$ أكبر من الواحد بقليل. قوة كل الوحدات المشغَّلة مجموع متوالية هندسية، ومعظمها يأتي من الوحدات الأخيرة. عندئذ تضيف الوحدة التالية المشغَّلة
\begin{empheq}[box=\keybox]{equation}
\frac{\Delta F}{F} \;\approx\; q-1 \;=\; \text{const} ,
\label{eq:sizeprinciple}
\end{empheq}
أي أن خطوة القوة تتناسب مع القوة نفسها. في الجهد الضعيف يجرعه الدماغ بحصص صغيرة، وفي القوي بحصص كبيرة. هذا \emph{عدد بفاصلة عائمة}: الأس يحدده عدد الوحدات المشغَّلة، والجزء العشري (المانتيسا) يحدده تردد نبضاتها. وهو المقياس اللوغاريتمي نفسه الذي لمستشعرات الفصل~\ref{sec:sensors}: مدخل الآلة ومخرجها يقيسان العالم بمقاييس نسبية.

وللفاصلة العائمة وجه آخر. فانتشار القوة يزداد هو أيضًا بما يتناسب معها: الحركة القوية أقل دقة حتمًا من الضعيفة. وبحسب فرضية مؤثرة، فهذا الضجيج المعتمد على الإشارة بالذات هو ما يفسر نعومة حركاتنا: الأمر الناعم يعطي أصغر انتشار في النقطة النهائية~\cite{HarrisWolpert1998}.

\section{تسحب ولا تدفع: سلكان لكل مفصل}

العضلة لا تستطيع إلا السحب. لا تستطيع الدفع، كما لا يستطيع عصبون \nt{GLUT} أن يُصدر إشارة سالبة. والحل معروف من الفصل~\ref{sec:glutcomputer}: \emph{سلكان}. كل مفصل يخدمه زوج من العضلات يسحبان في اتجاهين متعاكسين (الشكل~\ref{fig:antagonists}). إذا كانت قوتاهما $F_1$ و$F_2$، وذراع القوة $\ell$، فإن عزم الدوران وصلابة المفصل
\begin{empheq}[box=\keybox]{equation}
M \;=\; \ell\,(F_1-F_2), \qquad K \;\approx\; \kappa_m\,(F_1+F_2),
\label{eq:antagonists}
\end{empheq}
حيث $\kappa_m$ معامل يربط قوة العضلة بصلابتها: العضلة المشدودة أشد نابضية من المرتخية. الفرق بين القوتين يحدد العزم، ومجموعهما يحدد الصلابة~\cite{Hogan1984}. بسلكين يتحكم الدماغ في كميتين مستقلتين: إلى أين يُدار المفصل، وبأي إحكام يُمسك. وحين يلزم أن نمسك فنجانًا دون أن ينسكب، نشد العضلتين معًا: العزم نفسه، والصلابة أعلى، والدفعة العارضة تزيح اليد أقل.

\begin{figure}[t]
\centering
\begin{tikzpicture}[>=Stealth,
  nrn/.style={circle,draw,very thick,minimum size=0.95cm,inner sep=0pt,font=\scriptsize},
  inh/.style={-{Bar[width=7pt]},thick,red!70!black}]
% the joint: a bar on a pivot, two springs pulling down on either side
\fill[gray!30] (4.6,-0.25) -- (5.4,-0.25) -- (5,0.15) -- cycle;
\draw[line width=3pt,gray!50!black] (2.4,0.2) -- (7.6,0.2);
\fill[black] (5,0.2) circle (0.08);
\draw[decorate,decoration={coil,segment length=5pt,amplitude=4pt},very thick,blue!60!black] (3.0,0.2) -- (3.0,-1.6);
\draw[decorate,decoration={coil,segment length=5pt,amplitude=4pt},very thick,orange!85!black] (7.0,0.2) -- (7.0,-1.6);
\draw[very thick] (2.5,-1.6) -- (3.5,-1.6); \draw[very thick] (6.5,-1.6) -- (7.5,-1.6);
\node[font=\small,blue!60!black] at (2.35,-0.75) {$F_1$};
\node[font=\small,orange!85!black] at (7.65,-0.75) {$F_2$};
\draw[->,thick] (5.9,0.75) arc (60:120:1.8) node[above,font=\scriptsize,pos=0.5] {$M=\ell(F_1-F_2)$};
% motor neurons and reflex circuit
\node[nrn,fill=green!12] (M1) at (1.6,-3.0) {\nt{ACET}};
\node[nrn,fill=green!12] (M2) at (8.4,-3.0) {\nt{ACET}};
\node[nrn,fill=red!10] (G) at (5.0,-3.0) {\nt{GLYC}};
\node[nrn,fill=blue!6] (S) at (1.6,-5.0) {\nt{GLUT}};
\draw[->,very thick] (M1) -- (3.0,-1.7);
\draw[->,very thick] (M2) -- (7.0,-1.7);
\draw[->,thick] (S) -- (M1);
\draw[->,thick] (S) -- (G);
\draw[inh] (G) -- (M2);
\draw[->,thick,densely dashed,gray!60!black] (3.15,-1.2) to[bend left=25] node[pos=0.35,right,font=\scriptsize,black] {التمدد} (S.north east);
\draw[->,thick,blue!60!black] (-0.3,-3.0) node[left,font=\scriptsize,black,align=right] {أمر\\الدماغ} -- (M1);
\draw[->,thick,blue!60!black] (10.3,-3.0) node[right,font=\scriptsize,black,align=left] {أمر\\الدماغ} -- (M2);
\end{tikzpicture}
\caption{عضلتان لمفصل واحد، وأسرع حلقة تغذية راجعة. تتلقى عصبونات \nt{ACET} أوامر الدماغ وتسحب المفصل في اتجاهين متعاكسين؛ الفرق بين القوتين يديره، ومجموعهما يزيده صلابة. مستشعر التمدد (\nt{GLUT}) في العضلة اليسرى يثير عصبون \nt{ACET} الخاص بها، ويثبّط عبر وسيط \nt{GLYC} عصبون العضلة المضادة، وهذا هو المنع من الفصل~\ref{sec:gaba}.}
\label{fig:antagonists}
\end{figure}

\section{تغذية راجعة متأخرة}

لا تعمل العضلات عمياء. ففي كل منها مستشعرات تمدد، ذُكرت في الجدول~\ref{tab:sensors}، وأقصر حلقة تُغلق مباشرة في الحبل الشوكي (الشكل~\ref{fig:antagonists}). إذا مُدّت العضلة فجأة، أثار مستشعر التمدد عصبون \nt{ACET} الخاص بها، فتنقبض العضلة وتعيد المفصل. وفي الوقت نفسه تُطفأ العضلة المضادة عبر عصبون وسيط مثبِّط، كي لا تعيق. ويعمل وسيطًا هنا \nt{GLYC}، وهو النوع نفسه من الجدول~\ref{tab:types} الذي لم يُخصَّص له فصل: مكانه هنا بالذات، في حلقات الحبل الشوكي السريعة.

لكل حلقة تأخير $d$: للحلقة الشوكية نحو $25\text{--}30\ms$ في الذراع، وللحلقة عبر القشرة أطول بمرة ونصف إلى ثلاث مرات، وللحلقة عبر العين $100\text{--}200\ms$. والتغذية الراجعة المتأخرة، كما رأينا في القضية~\ref{prop:ei}، تُحدث تذبذبًا في النظام. ولأبسط منظِّم يصحح الانحراف $x$ بسرعة $G$ استنادًا إلى معلومات عمرها $d$ ثانية، $\dd x/\dd t=-G\,x(t-d)$، يكون شرط الاستقرار~\cite{Hayes1950}:
\begin{empheq}[box=\keybox]{equation}
G\,d \;<\; \frac{\pi}{2} ,
\label{eq:delaystable}
\end{empheq}
وعلى الحد يتذبذب النظام بتردد $1/(4d)$. اختبرنا ذلك على نموذج: عند $d=30\ms$ يخمد الانحراف عند $G=40$ في الثانية، أما عند $G=55$ في الثانية، فوق الحد $52$ بقليل، فيتذبذب متناميًا.

ومن هنا نتيجتان. الأولى: كلما طالت الحلقة وجب أن تصحح أبطأ. فعبر العين بتأخير مئة وخمسين ملّي ثانية لا يمكن تصحيح اليد أسرع من نحو عُشر ثانية، وإلا ارتجفت. والثانية: حد استقرار الحلقة الشوكية، $1/(4\cdot30\ms)\approx8$ ذبذبات في الثانية، قريب من تردد الرعشة الدقيقة الموجودة لدى كل إنسان سليم، أي من ثماني إلى اثنتي عشرة ذبذبة في الثانية. وحلقة التمدد أحد المصادر المحتملة لهذه الرعشة~\cite{McAuleyMarsden2000}.

فكيف يتحرك الدماغ إذن بسرعة ودقة مع أن التغذية الراجعة متأخرة؟ بحسب فرضية شائعة: \emph{بالتنبؤ}. فهو يحتفظ بنموذج داخلي للجسم ويحسب نتيجة الأمر دون انتظار المستشعرات~\cite{WolpertGhahramani2000}. والمستشعرات لازمة لتصحيح النموذج، لا لتصحيح كل حركة. وكان المهندس سيسمي هذا تحكمًا تنبؤيًا مع مراقِب للحالة.

\section{مولّدات إيقاع الحركات}

المشي والتنفس والمضغ حركات إيقاعية. فهل تحتاج إلى ساعة؟ لا. في الحبل الشوكي وجذع الدماغ شبكات صغيرة تولّد الإيقاع بنفسها، حتى لو لم يصلها أي شيء إيقاعي~\cite{MarderBucher2001}. وأبسط شبكة كهذه مجمَّعة من قطع لدينا من قبل: مجموعتان من عصبونات \nt{GLUT} تثبّط كل منهما الأخرى عبر وسطاء \nt{GABA} أو \nt{GLYC}، كما في الشكل~\ref{fig:wta}، وتتعبان ببطء، كذاكرة الفصل~\ref{sec:glutcomputer}:
\begin{empheq}[box=\keybox]{equation}
\tau\,\frac{\dd r_i}{\dd t} = -r_i + \bigl[I - \beta\,r_j - a_i\bigr]_+ ,
\qquad
\tau_a\,\frac{\dd a_i}{\dd t} = -a_i + b\,r_i ,
\label{eq:halfcentre}
\end{empheq}
حيث $a_i$ تعب المجموعة $i$، و$j$ المجموعة الأخرى. التثبيط المتبادل مع $\beta>1$ يختار فائزًا (القضية~\ref{prop:wta}). يعمل الفائز ويتعب؛ وحين يأكل التعب دخله، ينطلق الخاسر، وقد استراح، إلى المقدمة ويبدأ هو بالتعب. تعمل المجموعتان بالتناوب، وتتحكم كل منهما في عضلتها من الزوج (الشكل~\ref{fig:halfcentre}).

\begin{figure}[t]
\centering
\begin{tikzpicture}[>=Stealth,x=3.2cm,y=2.4cm]
\draw[->,gray!70!black] (0,0) -- (4.2,0) node[right,font=\small,black] {$t,\ \text{s}$};
\draw[->,gray!70!black] (0,0) -- (0,1.2) node[left,font=\small,black] {$r$};
\foreach \t in {0,1,2,3,4}{\draw[gray!60!black] (\t,-0.03) -- (\t,0.03); \node[below,font=\scriptsize] at (\t,0) {\t};}
\draw[thick,blue!60!black] plot file {figures/halfcentre1.dat};
\draw[thick,orange!85!black] plot file {figures/halfcentre2.dat};
\node[font=\scriptsize,blue!60!black,anchor=west] at (1.6,1.28) {المجموعة 1: ”إلى الأمام“};
\node[font=\scriptsize,orange!85!black,anchor=west] at (2.7,1.28) {المجموعة 2: ”إلى الخلف“};
\end{tikzpicture}
\caption{مولّد الإيقاع بحسب النموذج~\eqref{eq:halfcentre}: $I=1$، $\beta=2$، $b=2$، $\tau=20\ms$، $\tau_a=0.4\,\text{s}$. مجموعتان بتثبيط متبادل وتعب تعملان بالتناوب بدور $0.84\,\text{s}$، مع أن الدخل إشارة ثابتة.}
\label{fig:halfcentre}
\end{figure}

في نموذجنا، مع دخل ثابت، تتناوب المجموعتان بدور $0.84\,\text{s}$، أي نحو ضعف زمن التعب $\tau_a$. الأمر الثابت من الأعلى، ”امشِ“، يتحول إلى إيقاع في موضعه. وهذا إيقاع محلي آخر، كإيقاع حلقة \foreignlanguage{english}{\nt{GLUT}--\nt{GABA}} في الفصل~\ref{sec:gaba}: لا ينشأ إلا حين تعمل الشبكة، ولا يضبط الإيقاع إلا للعضلات التي تتحكم فيها.

\begin{keyblock}
\begin{proposition}[مخرج الآلة]
\label{prop:muscles}
يتحكم الدماغ في العضلات كهرمية من المنظِّمات. في الأعلى يُختار الفعل (الفصل~\ref{sec:dofa})؛ وأدنى منه تحوّل المولّدات المحلية الأوامرَ الثابتة إلى إيقاع، وتمسك الحلقاتُ السريعة في الحبل الشوكي المفاصلَ. وتُحدَّد القوة بفاصلة عائمة: الأس بعدد الوحدات المشغَّلة بحسب الحجم، والجزء العشري بتردد النبضات. وكل مفصل يتحكم فيه زوج من العضلات الساحبة، يحدد الفرق بين قوتيهما العزمَ، ومجموعهما الصلابةَ. هذه الأجهزة مقيسة؛ أما اعتماد الدماغ على نموذج داخلي للجسم ففرضية راسخة الأسس.
\end{proposition}
\end{keyblock}

\section{الخلاصة}

\begin{table}[t]
\centering
\footnotesize
\setlength{\tabcolsep}{4pt}
\renewcommand{\arraystretch}{1.15}
\begin{tabular}{>{\raggedright}p{3.0cm}>{\raggedright}p{4.6cm}>{\raggedright\arraybackslash}p{5.2cm}}
\toprule
ماذا يفعل المخرج & كيف & ما هو معروف \\
\midrule
يحوّل النبضات إلى قوة & مجموع النفضات~\eqref{eq:forcesum}، مرشح تمرير منخفض & مقيس؛ والجمع الخطي تبسيط \\
يجرّع القوة & تشغيل الوحدات بحسب الحجم، فاصلة عائمة~\eqref{eq:sizeprinciple} & ترتيب التشغيل مقيس \\
يدفع وهو لا يستطيع إلا السحب & زوج عضلات: العزم فرق، والصلابة مجموع~\eqref{eq:antagonists} & مقيس \\
يمسك المفصل & حلقة التمدد، منع المضادة عبر \nt{GLYC} & مقيس \\
لا يرتجف & $Gd<\pi/2$~\eqref{eq:delaystable} & ينتج من نظرية التحكم؛ والإسهام في الرعشة سبب محتمل \\
يتحرك بسرعة & التنبؤ بالنموذج الداخلي & فرضية راسخة الأسس \\
يمشي ويتنفس بإيقاع & مولّد الإيقاع~\eqref{eq:halfcentre} & المولّدات مقيسة؛ والمخطط الأبسط نموذج \\
\bottomrule
\end{tabular}
\caption{كيف تتحكم الآلة في العضلات.}
\label{tab:muscles}
\end{table}

مخرج الآلة مبني بالحيل نفسها التي بُني بها كل شيء آخر (الجدول~\ref{tab:muscles}): سلكان بدل العلامة، ومرشح تمرير منخفض بدل المحوِّل الرقمي التماثلي، ومقياس لوغاريتمي بدل الخطي، وتثبيط متبادل وتعب بدل مولّد نبضات الساعة. المدخل (الفصل~\ref{sec:sensors}) والمخرج يقيسان العالم بمقاييس نسبية، وبينهما تعمل الآلة التي خُصص لها هذا الكتاب.

\section{تمارين}

\begin{exercise}[\lvlA]
\label{ex:13:1}
\emph{الفاصلة العائمة بالأرقام.} في العضلة $N=100$ وحدة حركية قواها $f_n=f_1q^{\,n-1}$؛ أقواها أقوى من أضعفها $100$ مرة. (أ)~أوجد $q$، والخطوة النسبية~\eqref{eq:sizeprinciple}، والمدى الديناميكي بالديسيبل. (ب)~ما الخطوة عند القوة $10f_1$ في ”محوِّل رقمي تماثلي خطي“ من $100$ وحدة متماثلة بالقوة الكلية نفسها؟
\end{exercise}

\begin{exercise}[\lvlA]
\label{ex:13:2}
\emph{ثمن هامش الأمان.} مع كل نبضة يطلق المشبك على الليف العضلي عددًا بواسونيًا من الحزم متوسطه $m$، والليف يطلق عند $n_{\text{th}}=20$ حزمة فأكثر؛ وهامش الأمان $S=m/n_{\text{th}}$. كم إخفاقًا في اليوم تعطي وحدة تعمل بتردد $20$ نبضة في الثانية عند $S=2$ وعند $S=4$؟
\end{exercise}

\begin{exercise}[\lvlB]
\label{ex:13:3}
\emph{العضلة مرشحًا.} النفضة~\eqref{eq:twitch} مع $T_c=50\ms$ هي الاستجابة النبضية لنظام خطي. (أ)~أوجد دالة تحويله $H(s)$ وتردد القطع. (ب)~عند أي تردد للنبضات المنتظمة يبلغ مدى تموجات القوة $10\%$ من متوسط القوة؟
\end{exercise}

\begin{exercise}[\lvlB]
\label{ex:13:4}
\emph{لا تصحيح أسرع من التأخير.} المنظِّم $\dd x/\dd t=-G\,x(t-d)$. (أ)~بيّن أن الانحراف يخمد بلا تذبذب بأسرع ما يمكن عند $Gd=1/e$، بمعدل $1/d$. (ب)~أوجد هذا $G$ وثابت زمن الخمود للحلقة الشوكية، $d=30\ms$، وللحلقة عبر العين، $d=150\ms$.
\end{exercise}

\begin{exercise}[\lvlB]
\label{ex:13:5}
\emph{دور مولّد الإيقاع.} عند $\tau\ll\tau_a$ يُحدَّد الدور $T$ للنموذج~\eqref{eq:halfcentre} بالمعادلة $(\beta-1)(1-E^{2+b})/(\beta-E)=b\,(1-E^{1+b})/(1+b)$، حيث $E=e^{-T/(2\tau_a)}$. (أ)~أوجد $T$ عند $\beta=b=2$ و$\tau_a=0.4\,\text{s}$. (ب)~بيّن أن $T$ لا يتوقف على الدخل $I$، وأن الإيقاع ممكن فقط عند $b>\beta-1$.
\end{exercise}

\begin{exercise}[\lvlB]
\label{ex:13:6}
\emph{نمذجة المولّد.} استورد الدالة \foreignlanguage{english}{\texttt{halfcentre}} من \foreignlanguage{english}{\texttt{models/\allowbreak muscle\_models.py}} (لا تشغّل الملف نفسه)، وقس الدور، مهملًا الدورتين الأوليين، عند $b=0.8$، $1.05$، $1.2$، $1.5$، $2$، $4$ وعند $I=0.5$، $1$، $2$. هل يستطيع الأمر ”امشِ أسرع“ أن يغيّر الوتيرة عبر $I$؟
\end{exercise}

\begin{exercise}[\lvlB]
\label{ex:13:7}
\emph{الصلابة في مواجهة المنعكس.} يحيط بالمفصل حلقة تمدد: $\dd\theta/\dd t=-a\theta(t)-G\theta(t-d)$، $a=K/B$، $d=30\ms$. ردّها إلى التمرين~\ref{ex:13:4}، وأوجد أصغر $a$ يخمد عنده الانحراف بلا تذبذب بثابت زمن $15\ms$، و$G$ اللازم. كيف يجب تغيير $G$ عند زيادة الشد المشترك للعضلات؟
\end{exercise}

\begin{exercise}[\lvlC]
\label{ex:13:8}
\emph{مقبض الوتيرة.} بحسب التمرينين~\ref{ex:13:5} و\ref{ex:13:6} لا يغيّر الدخل $I$ للمولّد~\eqref{eq:halfcentre} إلا السعة، لا الوتيرة. اقترح أدنى تعديل للنموذج من قطع الكتاب، بحيث لا يوجد إيقاع دون $I_{\min}$ ما، ويتناقص الدور فوقها مع ازدياد $I$، إلى النصف على الأقل عند مضاعفة $I$ ثلاث مرات. أوجد $I_{\min}$ عند $\tau\ll\tau_a$ وتحقق بالنمذجة.
\end{exercise}
